%%%%%%%%%%%%%%%%%%%%%%%%%%%%%%%%%%%%%%%%%%%%%%%%%%%%%%%%%%%%%%%%%%%%%%%%%%%%%%%%
% EMCH 367 / ELCT 331 - Modern Control Systems for Physical Engineering Applications
% University of South Carolina - Fall 2026
%%%%%%%%%%%%%%%%%%%%%%%%%%%%%%%%%%%%%%%%%%%%%%%%%%%%%%%%%%%%%%%%%%%%%%%%%%%%%%%%

\documentclass[11pt,letterpaper]{article}

%-------------------------------------------------------------------------------
% PACKAGES
%-------------------------------------------------------------------------------
\usepackage[margin=1in]{geometry}
\usepackage[dvipsnames,table]{xcolor}
\usepackage{titlesec}
\usepackage{booktabs}
\usepackage{enumitem}
\usepackage{hyperref}
\usepackage{fancyhdr}
\usepackage{tabularx}
\usepackage{array}
\usepackage{multirow}
\usepackage{graphicx}
\usepackage{parskip}
\usepackage{longtable}
\usepackage{colortbl}
\usepackage{tabularx}
\usepackage{tabularray}
\UseTblrLibrary{booktabs}

%-------------------------------------------------------------------------------
% COLOR DEFINITIONS - UofSC Brand Colors
%-------------------------------------------------------------------------------
\definecolor{UofSCGarnet}{HTML}{73000a}
\definecolor{UofSCBlack}{HTML}{000000}
\definecolor{UofSCWhite}{HTML}{ffffff}
\definecolor{UofSCAtlantic}{HTML}{466A9F}
\definecolor{garnet}{HTML}{73000a}
\definecolor{atlantic}{HTML}{466A9F}
\definecolor{congaree}{HTML}{1F414D}
\definecolor{examrow}{HTML}{FFF2E3}
\definecolor{presentrow}{HTML}{E8F0F8}
\definecolor{breakrow}{HTML}{FCEAEA}
\definecolor{bufferrow}{HTML}{E8F5E0}

%-------------------------------------------------------------------------------
% HYPERREF SETUP
%-------------------------------------------------------------------------------
\hypersetup{
    colorlinks=true,
    linkcolor=UofSCGarnet,
    urlcolor=UofSCAtlantic,
    citecolor=UofSCGarnet,
    pdfauthor={J.C. Vaught},
    pdftitle={EMCH 367 / ELCT 331 Syllabus - Fall 2026},
    pdfsubject={Modern Control Systems for Physical Engineering Applications}
}

%-------------------------------------------------------------------------------
% SECTION FORMATTING
%-------------------------------------------------------------------------------
\titleformat{\section}
    {\Large\bfseries\color{UofSCGarnet}}
    {\thesection}{1em}{}[\vspace{-1.2em}\textcolor{UofSCGarnet}{\rule{\textwidth}{1pt}}]

\titleformat{\subsection}
    {\large\bfseries\color{UofSCBlack}}
    {\thesubsection}{1em}{}

\titlespacing*{\section}{0pt}{1.5ex plus 1ex minus .2ex}{1ex plus .2ex}
\titlespacing*{\subsection}{0pt}{1ex plus 1ex minus .2ex}{0.5ex plus .2ex}

%-------------------------------------------------------------------------------
% HEADER/FOOTER
%-------------------------------------------------------------------------------
\pagestyle{fancy}
\fancyhf{}
\renewcommand{\headrulewidth}{0.5pt}
\renewcommand{\headrule}{\hbox to\headwidth{\color{UofSCGarnet}\leaders\hrule height \headrulewidth\hfill}}
\fancyhead[L]{\textcolor{UofSCGarnet}{\textbf{EMCH 367 / ELCT 331}}}
\fancyhead[R]{\textcolor{UofSCGarnet}{Fall 2026}}
\fancyfoot[C]{\thepage}

%-------------------------------------------------------------------------------
% CUSTOM COMMANDS
%-------------------------------------------------------------------------------
\newcommand{\courseheader}{%
    \begin{center}
        {\color{UofSCGarnet}\rule{\textwidth}{2pt}}\\[0.5em]
        {\Huge\bfseries\color{UofSCGarnet} EMCH 367 / ELCT 331}\\[0.3em]
        {\LARGE\bfseries Modern Control Systems for}\\[0.1em]
        {\LARGE\bfseries Physical Engineering Applications}\\[0.5em]
        {\large University of South Carolina}\\[0.3em]
        {\large\textbf{Fall 2026}}\\[0.5em]
        {\color{UofSCGarnet}\rule{\textwidth}{2pt}}
    \end{center}
}

\newcommand{\infobox}[2]{%
    \textcolor{UofSCGarnet}{\textbf{#1:}} #2
}

%-------------------------------------------------------------------------------
% TABLE COLUMN TYPES
%-------------------------------------------------------------------------------
\newcolumntype{L}[1]{>{\raggedright\arraybackslash}p{#1}}
\newcolumntype{C}[1]{>{\centering\arraybackslash}p{#1}}
\newcolumntype{R}[1]{>{\raggedleft\arraybackslash}p{#1}}

%%%%%%%%%%%%%%%%%%%%%%%%%%%%%%%%%%%%%%%%%%%%%%%%%%%%%%%%%%%%%%%%%%%%%%%%%%%%%%%%
\begin{document}
\thispagestyle{empty}

\courseheader

\vspace{1em}

%-------------------------------------------------------------------------------
% COURSE INFORMATION BOX
%-------------------------------------------------------------------------------
\noindent
\fcolorbox{UofSCGarnet}{UofSCWhite}{%
\begin{minipage}[t]{0.48\textwidth}
    \vspace{0.5em}
    \infobox{Course ID}{EMCH 367 / ELCT 331}\\[0.5em]
    \infobox{Meeting Time}{Tue/Thu 10:05--11:20 AM}\\[0.5em]
    \infobox{Location}{300MN B112}
    \vspace{0.5em}
\end{minipage}%
}%
\hfill%
\fcolorbox{UofSCGarnet}{UofSCWhite}{%
\begin{minipage}[t]{0.48\textwidth}
    \vspace{0.5em}
    \infobox{Instructor}{J.C. Vaught}\\[0.5em]
    \infobox{Office}{HRZN 011; T/TH 11:20--12:30}\\[0.5em]
    \infobox{Reading}{Textbook TBD}
    \vspace{0.5em}
\end{minipage}%
}

\vspace{1.5em}

%-------------------------------------------------------------------------------
\section*{Catalog Description}
%-------------------------------------------------------------------------------

This course develops practical competency in modeling, feedback control, estimation, and implementation for real engineered systems. The course is concept-first in lecture and application-heavy in assignments. Students learn to build control-ready models of physical systems, design and evaluate feedback controllers, implement controllers digitally, handle key nonlinear effects, and work with multivariable systems using state space methods. The course includes two team projects with physical demonstrations and presentations.

%-------------------------------------------------------------------------------
\section*{Prerequisites}
%-------------------------------------------------------------------------------

Students are expected to be fluent in calculus through multivariable calculus, differential equations, and linear algebra. Students are expected to be comfortable with at least one programming language that can produce plots. The course does not include a formal review of math prerequisites.

%-------------------------------------------------------------------------------
\section*{Required Resources}
%-------------------------------------------------------------------------------

Students need a laptop capable of running the course programming environment. The course will standardize on a primary environment (Python or MATLAB) and provide a reference environment for grading. Students must be able to run provided starter code, produce required plots, and submit reproducible results. Hardware resources and safety requirements for demos are provided by the course and described on the course site.

%-------------------------------------------------------------------------------
\section*{Instructional Approach}
%-------------------------------------------------------------------------------

Lecture emphasizes concepts, design reasoning, and engineering workflows rather than step-by-step derivations. Mathematical derivations and deeper proofs are provided through optional typed notes and video resources. Theory homework requires formal engagement with selected mathematical tools, while application homework emphasizes implementation, validation, and interpretation.

%-------------------------------------------------------------------------------
\section*{Assessment Weights}
%-------------------------------------------------------------------------------

This course is not curved. Final grades are determined by the following components.

\vspace{1em}
\noindent
\begin{tabularx}{\linewidth}{X >{\raggedleft\arraybackslash}p{2cm}}
\toprule
\textbf{Component} & \textbf{Weight} \\
\midrule
Homework (11 sets, lowest dropped) & 50\% \\
Projects & 40\% \\
\quad Project 1 & 15\% \\
\quad Project 2 & 25\% \\
Exams (two in-term, 5\% each) & 10\% \\
\bottomrule
\end{tabularx}

%-------------------------------------------------------------------------------
\section*{Grading Scale}
%-------------------------------------------------------------------------------

Final letter grades are assigned according to the following scale.

\vspace{0.5em}
\noindent
\renewcommand{\arraystretch}{1.2}
\begin{tabularx}{\linewidth}{>{\bfseries}p{2.2cm} *{8}{>{\centering\arraybackslash}X}}
\toprule
\textbf{Grade} & A & B+ & B & C+ & C & D+ & D & F \\
\midrule
\textbf{Percentage} & 90--100 & 85--89 & 80--84 & 75--79 & 70--74 & 65--69 & 60--64 & $<$60 \\
\bottomrule
\end{tabularx}

%-------------------------------------------------------------------------------
\section*{Homework}
%-------------------------------------------------------------------------------

There are 11 homework sets total. The lowest score is dropped. The remaining 10 are weighted by rank, with the lowest grade receiving the highest weight.

\vspace{0.5em}
\noindent
\renewcommand{\arraystretch}{1.2}
\begin{tabularx}{\linewidth}{X *{10}{>{\centering\arraybackslash}X}}
\toprule
\textbf{Rank} & 1st & 2nd & 3rd & 4th & 5th & 6th & 7th & 8th & 9th & 10th \\
\midrule
\textbf{Weight} & 10\% & 8\% & 6\% & 6\% & 4\% & 4\% & 3\% & 3\% & 3\% & 3\% \\
\bottomrule
\end{tabularx}
\vspace{0.5em}

Each homework assignment consists of two individual assignments: one for coding/computer submission and one for theory derivation with math.

%-------------------------------------------------------------------------------
\section*{Submission Requirements}
%-------------------------------------------------------------------------------

Each homework is due Monday at midnight before we cover another section. Assignments are required to be typeset (e.g., \LaTeX, Word with equation editor) and any handwritten submissions will receive a zero. All assignments are graded on final correctness of each part on a pass/fail basis.

%-------------------------------------------------------------------------------
\section*{Exams}
%-------------------------------------------------------------------------------

There will be two required exams, each worth 5\%, as well as an optional final exam. The exams will focus on modeling assumptions, control structure, interpreting behavior, engineering steps, and diagnosing failures. All exams will contain up to 10\% bonus. The final exam is cumulative and will replace either a homework or exam grade if it improves your grade; otherwise it will not be applied.

%-------------------------------------------------------------------------------
\section*{Additional Credit}
%-------------------------------------------------------------------------------

Students may earn additional credit of up to 5\% per assignment or project by demonstrating work that exceeds expectations. An optional research project may also be approved until two weeks before the final day of class. A 5-point research project is expected to involve a focused investigation of a specific control technique or application, with a brief written summary and demonstration. A 10-point research project is expected to involve a more comprehensive study, including implementation, experimentation, and a detailed written report with analysis of results.

%-------------------------------------------------------------------------------
\section*{Learning Outcomes}
%-------------------------------------------------------------------------------

By the end of the course, students will be able to achieve the following objectives.

\begin{table}[htbp]
\begin{tabularx}{\linewidth}{c L{3.5cm} X}
\toprule
\textbf{\#} & \textbf{Outcome} & \textbf{Description} \\
\midrule
1 & Model physical systems & Translate a physical system into a control-ready model with explicit assumptions and validity ranges, including actuator and sensor nonidealities \\
\addlinespace
2 & Design feedback controllers & Design baseline feedback controllers and select appropriate control architectures for real systems \\
\addlinespace
3 & Evaluate performance & Evaluate performance using standard time-domain metrics and interpret stability evidence in both time and frequency perspectives \\
\addlinespace
4 & Implement digitally & Implement controllers in discrete time and anticipate common sim-to-hardware failure modes \\
\addlinespace
5 & Handle nonlinearities & Model and reason about nonlinear effects and robustness \\
\addlinespace
6 & Work with multivariable systems & Represent and reason about multivariable systems in state space and understand the motivation for native multivariable controller designs \\
\addlinespace
7 & Communicate effectively & Communicate engineering decisions and experimental evidence in written and oral form \\
\bottomrule
\end{tabularx}
\end{table}

%-------------------------------------------------------------------------------
\section*{Course Policies}
%-------------------------------------------------------------------------------

%-------------------------------------------------------------------------------
\subsection*{AI Tools Policy}
%-------------------------------------------------------------------------------

Artificial intelligence tools are allowed and highly encouraged in this course. Given the complexity of modern control system design and implementation, students who do not utilize AI assistance are likely to struggle significantly. AI may be used without restriction on all homework assignments, projects, and other coursework. Disclosure of AI use is not required and is discouraged. The only exception is examinations, where AI tools are not permitted.

%-------------------------------------------------------------------------------
\subsection*{Collaboration}
%-------------------------------------------------------------------------------

Conceptual discussion with classmates is allowed and encouraged, but all submissions must be your own work. Project work must reflect equitable contribution from all team members. External resources must be cited.

%-------------------------------------------------------------------------------
\subsection*{Late Work}
%-------------------------------------------------------------------------------

Late homework submissions are generally not accepted because one lowest score is dropped. Project deadlines are firm due to hardware scheduling and presentation logistics.

%-------------------------------------------------------------------------------
\subsection*{Accessibility}
%-------------------------------------------------------------------------------

Students with accommodations should coordinate with the instructor early in the semester to ensure appropriate arrangements can be made.

%%%%%%%%%%%%%%%%%%%%%%%%%%%%%%%%%%%%%%%%%%%%%%%%%%%%%%%%%%%%%%%%%%%%%%%%%%%%%%%
% TEAM MEMBERS & SELECTION SECTION
%%%%%%%%%%%%%%%%%%%%%%%%%%%%%%%%%%%%%%%%%%%%%%%%%%%%%%%%%%%%%%%%%%%%%%%%%%%%%%%

\section*{Project Team Members \& Selection}
\label{sec:team-formation}

Projects constitute a significant portion of the course grade and provide hands-on experience with control system design, implementation, and validation. Students will complete two projects throughout the semester, progressing from single-input single-output (SISO) systems to more complex multi-input multi-output (MIMO) systems with machine learning integration.

All projects must be completed in teams of 2--3 students. Single-person projects are not permitted due to insufficient time for individual presentations and the collaborative nature of real-world engineering practice.

Students may self-organize into teams during the first lecture. After the first lecture, any remaining unassigned students will be placed into teams by the instructor. No team changes are permitted until after Project~1 is complete. Requests for changes after Project~1 must be submitted in writing.

%%%%%%%%%%%%%%%%%%%%%%%%%%%%%%%%%%%%%%%%%%%%%%%%%%%%%%%%%%%%%%%%%%%%%%%%%%%%%%%
% PROJECT 1 SECTION
%%%%%%%%%%%%%%%%%%%%%%%%%%%%%%%%%%%%%%%%%%%%%%%%%%%%%%%%%%%%%%%%%%%%%%%%%%%%%%%

\section*{Project 1: SISO Control Implementation}
\label{sec:project1}

Project~1 focuses on the design and implementation of a controller for a single-input single-output (SISO) system. This project establishes foundational skills in control implementation that will be extended in Project~2.

Systems may be continuous-time or discrete-time. Implementation must be on real hardware or a detailed simulation with realistic sensor and actuator models (including noise, quantization, delays, and saturation). Projects are open-ended---students are encouraged to propose their own systems.

The following systems are provided as a starting menu. Students may select from this list or propose an alternative system for approval.

\begin{table}[htbp]
\begin{tabularx}{\linewidth}{lXX}
\toprule
\textbf{System} & \textbf{Input} & \textbf{Output} \\
\midrule
Inverted pendulum / cart-pole & Cart force or motor voltage & Pendulum angle \\
Vehicle steering control & Steering angle & Lateral position or heading \\
DC motor position control & Motor voltage & Shaft position \\
DC motor speed control & Motor voltage & Shaft angular velocity \\
Ball and beam & Beam angle & Ball position \\
Thermal system regulation & Heater power & Temperature \\
\textit{Custom approved system} & \textit{TBD} & \textit{TBD} \\
\bottomrule
\end{tabularx}
\end{table}

%%%%%%%%%%%%%%%%%%%%%%%%%%%%%%%%%%%%%%%%%%%%%%%%%%%%%%%%%%%%%%%%%%%%%%%%%%%%%%%
% PROJECT 2 SECTION
%%%%%%%%%%%%%%%%%%%%%%%%%%%%%%%%%%%%%%%%%%%%%%%%%%%%%%%%%%%%%%%%%%%%%%%%%%%%%%%

\section*{Project 2: MIMO with ML-Based Control}
\label{sec:project2}

Project~2 advances to multi-input multi-output (MIMO) systems and introduces machine learning-based control methods. This project emphasizes the comparison between classical and modern control approaches.

Systems should have multiple inputs and outputs (preferably many I/O channels). Implementation must be discrete-time. Each project must include both an ML-based controller (e.g., reinforcement learning, neural network control, adaptive methods) and a classical controller for comparison (e.g., LQR, PID, state feedback). Projects are open-ended---students are encouraged to continue/expand their Project~1 system or propose a new system.

\begin{table}[htbp]
\begin{tabularx}{\linewidth}{lXX}
\toprule
\textbf{System} & \textbf{Inputs} & \textbf{Outputs} \\
\midrule
Quadrotor attitude/position & Motor speeds (4) & Roll, pitch, yaw, position (6) \\
Multi-joint robotic arm & Joint torques/voltages & Joint angles, end-effector pose \\
Multi-zone HVAC system & Damper positions, fan speeds & Zone temperatures, humidity \\
Autonomous vehicle & Steering, throttle, braking & Position, velocity, heading \\
\textit{Custom approved system} & \textit{TBD} & \textit{TBD} \\
\bottomrule
\end{tabularx}
\end{table}

%%%%%%%%%%%%%%%%%%%%%%%%%%%%%%%%%%%%%%%%%%%%%%%%%%%%%%%%%%%%%%%%%%%%%%%%%%%%%%%
% PROJECT RESOURCES SECTION
%%%%%%%%%%%%%%%%%%%%%%%%%%%%%%%%%%%%%%%%%%%%%%%%%%%%%%%%%%%%%%%%%%%%%%%%%%%%%%%

\section*{Project Resources}
\label{sec:resources}

The following resources are available to all project teams.

\begin{table}[htbp]
\begin{tabularx}{\linewidth}{lX}
\toprule
\textbf{Resource} & \textbf{Details} \\
\midrule
3D Printers & Available for fabricating custom mechanical parts \\
Computers & Workstations with full course software environment \\
Arduino Microcontrollers & Standard issue for all teams \\
Raspberry Pi & Available upon request \\
Sensor/Actuator Kits & Standardized kits provided to each team \\
Personal Equipment & Students may purchase additional equipment at their own expense \\
Research Lab Equipment & Students may integrate equipment from their research laboratories with advisor approval \\
\bottomrule
\end{tabularx}
\end{table}

%%%%%%%%%%%%%%%%%%%%%%%%%%%%%%%%%%%%%%%%%%%%%%%%%%%%%%%%%%%%%%%%%%%%%%%%%%%%%%%
% PROJECT DEMONSTRATIONS SECTION
%%%%%%%%%%%%%%%%%%%%%%%%%%%%%%%%%%%%%%%%%%%%%%%%%%%%%%%%%%%%%%%%%%%%%%%%%%%%%%%

\section*{Project Demonstrations}
\label{sec:demo-requirements}

All projects require demonstration of the implemented control system. A recorded video demonstration is required and must be uploaded with the project submission. Live demonstrations during presentations are permitted only if the equipment is small and portable; large or fragile setups should be demonstrated via video only.

Demonstrations must show multiple trials to establish repeatability of performance---a single successful run is insufficient. Demonstrations must also include failure case analysis: What conditions cause the controller to fail? What are the stability boundaries? How does the system degrade under disturbances or parameter variations?

\newpage

%%%%%%%%%%%%%%%%%%%%%%%%%%%%%%%%%%%%%%%%%%%%%%%%%%%%%%%%%%%%%%%%%%%%%%%%%%%%%%%
% PROJECT DELIVERABLES SECTION
%%%%%%%%%%%%%%%%%%%%%%%%%%%%%%%%%%%%%%%%%%%%%%%%%%%%%%%%%%%%%%%%%%%%%%%%%%%%%%%

\section*{Project Deliverables}
\label{sec:deliverables}

Both projects require the following deliverables. All items are mandatory unless explicitly stated otherwise.

\begin{table}[htbp]
\begin{tabularx}{\linewidth}{clX}
\toprule
\textbf{\#} & \textbf{Deliverable} & \textbf{Description} \\
\midrule
1 & Project Proposal Document & Submitted for instructor approval \textit{before} work begins. Must include system description, objectives, methodology, and timeline. \\
\addlinespace
2 & Code Repository & All source code must be submitted. Code must be reproducible with clear instructions for execution. \\
\addlinespace
3 & Team Report & Comprehensive technical report covering theory, design, implementation, results, and analysis. \\
\addlinespace
4 & Individual Report & Each team member submits a personal reflection on their contributions and key learnings. \\
\addlinespace
5 & Presentation Slides & Submitted \textit{before} presentation day for review. \\
\addlinespace
6 & Live Presentation & In-class presentation followed by Q\&A session with instructor and peers. \\
\addlinespace
7 & Demo Video & Recorded demonstration with narration explaining the system operation and results. \\
\bottomrule
\end{tabularx}
\end{table}

\newpage

%%%%%%%%%%%%%%%%%%%%%%%%%%%%%%%%%%%%%%%%%%%%%%%%%%%%%%%%%%%%%%%%%%%%%%%%%%%%%%%
% ASSIGNMENT SUMMARY SECTION
%%%%%%%%%%%%%%%%%%%%%%%%%%%%%%%%%%%%%%%%%%%%%%%%%%%%%%%%%%%%%%%%%%%%%%%%%%%%%%%

\section*{Assignment Summary}

Each homework (HW1--HW10) consists of two files: HW\_X\textsubscript{m} (math) and HW\_X\textsubscript{c} (code).

\begin{table}[h]
\begin{tabularx}{\linewidth}{llX}
\toprule
\textbf{Assignment} & \textbf{Due Date} & \textbf{Topic Coverage} \\
\midrule
HW0 & Sun 8/23 (midnight) & Mathematics Prerequisites \\
HW1 & Mon 8/24 & Modeling I (Lecture 1) \\
HW2 & Mon 8/31 & Modeling II: Actuators \& Sensors (Lectures 2--3) \\
HW3 & Mon 9/7 & Basic Feedback Design (Lectures 4--5) \\
HW4 & Mon 9/14 & Stability Analysis (Lectures 6--7) \\
HW5 & Mon 9/21 & Architecture \& State Space (Lectures 8--9) \\
HW6 & Mon 10/5 & Discrete \& Nonlinear (Lectures 11--13) \\
HW7 & Mon 10/12 & State Estimation (Lecture 13) \\
HW8 & Mon 10/26 & MIMO Design (Lectures 16--18) \\
HW9 & Mon 11/9 & Learned Perception (Lectures 21--22) \\
HW10 & Mon 11/23 & Learned Control (Lectures 23--26) \\
\midrule
\textbf{Exam 1} & Tue 9/22 & Modeling, Feedback, Stability \\
\textbf{Exam 2} & Thu 10/29 & MIMO, Robustness, Implementation \\
\midrule
Project 1 Proposal & Thu 9/24 & --- \\
Project 1 Report & Thu 10/8 & --- \\
Project 2 Proposal & Tue 11/3 & --- \\
Project 2 Report & Thu 12/3 & --- \\
\bottomrule
\end{tabularx}
\label{tab:assignments}
\end{table}

\newpage

%%%%%%%%%%%%%%%%%%%%%%%%%%%%%%%%%%%%%%%%%%%%%%%%%%%%%%%%%%%%%%%%%%%%%%%%%%%%%%%
% SCHEDULE SECTION
%%%%%%%%%%%%%%%%%%%%%%%%%%%%%%%%%%%%%%%%%%%%%%%%%%%%%%%%%%%%%%%%%%%%%%%%%%%%%%%

\section*{Course Schedule}

The course meets Tuesdays and Thursdays throughout Fall 2026. The schedule below maps all 28 meetings to specific dates, accounting for university holidays and breaks.

\vspace{1em}

\begin{longtable}{c c l p{10cm}}
\toprule
\textbf{Date} & \textbf{Lec} & \textbf{Type} & \textbf{Topic} \\
\midrule
\endfirsthead

\toprule
\textbf{Date} & \textbf{Lec} & \textbf{Type} & \textbf{Topic} \\
\midrule
\endhead

\midrule
\multicolumn{4}{r}{\textit{Continued on next page}} \\
\endfoot

\bottomrule
\endlastfoot

% Week 1
Thu 8/20 & 1 & Content & Modeling Physical Systems for Control \\

% Week 2
Tue 8/25 & 2 & Content & Modeling Physical Actuators for Control \\
Thu 8/27 & 3 & Content & Modeling Sensor Error for Control \\

% Week 3
Tue 9/1 & 4 & Content & Designing Control Systems \\
Thu 9/3 & 5 & Content & Measuring Controller Performance \\

% Week 4 (Labor Day Monday)
Tue 9/8 & 6 & Content & Stability Analysis in Time and Frequency Domain \\
Thu 9/10 & 7 & Content & Control Architecture \\

% Week 5
Tue 9/15 & 8 & Content & Introducing State Space Representation \\
\rowcolor{bufferrow}
Thu 9/17 & 9 & Buffer & Integration Studio and Exam 1 Prep \\

% Week 6 - Exam 1
\rowcolor{examrow}
Tue 9/22 & 10 & Exam & \textbf{Exam 1}: Modeling, Feedback, Stability \\
Thu 9/24 & 11 & Content & Implementing Control in Discrete Domain \\

% Week 7
Tue 9/29 & 12 & Content & Nonlinear System Accounting and Control Robustness \\
Thu 10/1 & 13 & Content & Estimating State in Nonlinear Systems \\

% Week 8 - Project 1 Presentations
\rowcolor{presentrow}
Tue 10/6 & 14 & Presentation & Project 1 Presentations Day 1 \\
\rowcolor{presentrow}
Thu 10/8 & 15 & Presentation & Project 1 Presentations Day 2 \\

% Week 9 (Fall Break Thu)
Tue 10/13 & 16 & Content & Modeling MIMO Systems in State Space \\
\rowcolor{breakrow}
Thu 10/15 & --- & \textit{No Class} & \textit{Fall Break} \\

% Week 10
Tue 10/20 & 17 & Content & Designing Control Systems for MIMO \\
Thu 10/22 & 18 & Content & Native MIMO Controllers: LQR, LQG, MPC \\

% Week 11 - Buffer + Exam 2
\rowcolor{bufferrow}
Tue 10/27 & 19 & Buffer & Integration Studio and Exam 2 Prep \\
\rowcolor{examrow}
Thu 10/29 & 20 & Exam & \textbf{Exam 2}: MIMO, Robustness, Implementation \\

% Week 12 (ML-based control begins)
Tue 11/3 & 21 & Content & CNN and Transformer Encoders for State Estimation \\
Thu 11/5 & 22 & Content & Transformer World Models for Dynamics \\

% Week 13
Tue 11/10 & 23 & Content & Hybrid Model-Based and Learned Control \\
Thu 11/12 & 24 & Content & Differentiable Predictive Control \\

% Week 14
Tue 11/17 & 25 & Content & Diffusion Policies for Trajectory Tracking \\
Thu 11/19 & 26 & Content & Generalization and Sim-to-Real \\

% Week 15 - Thanksgiving
\rowcolor{breakrow}
Tue 11/24 & --- & \textit{No Class} & \textit{Thanksgiving Break} \\
\rowcolor{breakrow}
Thu 11/26 & --- & \textit{No Class} & \textit{Thanksgiving Break} \\

% Week 16 - Final Presentations
\rowcolor{presentrow}
Tue 12/1 & 27 & Presentation & Final Project Presentations Day 1 \\
\rowcolor{presentrow}
Thu 12/3 & 28 & Presentation & Final Project Presentations Day 2 \\

\end{longtable}

\end{document}
